\documentclass[11pt]{article}
% Standalone preamble for web-only supplementary notes.
\usepackage[T1]{fontenc}
\usepackage{amsmath,amssymb,amsthm}
\usepackage{hyperref}
\usepackage{cleveref}
\DeclareMathOperator{\Map}{Map}
\DeclareMathOperator{\Fun}{Fun}
\newcommand{\iso}{\xrightarrow{\sim}}
\newtheorem{lemma}{Supplementary Lemma}
\newtheorem{corollary}[lemma]{Supplementary Corollary}
\crefname{lemma}{Supplementary Lemma}{Supplementary Lemmas}
\Crefname{lemma}{Supplementary Lemma}{Supplementary Lemmas}
\crefname{corollary}{Supplementary Corollary}{Supplementary Corollaries}
\Crefname{corollary}{Supplementary Corollary}{Supplementary Corollaries}

\begin{document}
\section*{Restricting and pasting cone diagrams}

The \href{pullbacks.html\#lem:Pasting_Lemma_Pullbacks}{pullback pasting lemma}
uses compatibility of chosen factorizations with restriction and pasting.
These calculations must retain the matching identification of each cone.
All categories below are absolute, with the Chapter 1 coherence and
pullback structure.

For $f\colon C\to E$ and $g\colon D\to E$, a cone on $T$ is
$s=(p,q,\tau)$ with $\tau\colon fp=gq$.
A comparison to $s'=(p',q',\tau')$ consists of $L\colon p=p'$,
$R\colon q=q'$, and
\[
 \tau'\cdot(f\star L)=(g\star R)\cdot\tau.
\]
The symbol $\star$ denotes whiskering; vertical composition applies the
rightmost factor first. Write $\alpha_{r,p,f}\colon(fp)r=f(pr)$.
These descriptions concern data of cones and comparisons, without requiring
a category of all cones.

\begin{lemma}
\label[lemma]{lem:Supplementary_Cone_Restriction}
Restriction by $r\colon S\to T$ sends $s$ to the cone $r^*s$ with
legs $pr,qr$ and matching
\[
 \alpha_{r,q,g}\cdot(\tau\star r)\cdot\alpha_{r,p,f}^{-1}.
\]
It takes cone comparisons to cone comparisons. There is a comparison
$\mathrm{id}^*s\to s$ with the right unitors as leg comparisons.
For $u\colon U\to T$ and $r\colon S\to U$, there is also a comparison
$r^*(u^*s)\to(u\circ r)^*s$ with the associators as leg comparisons.
\end{lemma}
\begin{proof}[Proof idea]
Prewhisker the two leg comparisons and their matching equation.
The mixed whiskering law and associator compatibility put the resulting
equation into the displayed form. For identity restriction, use
compatibility of the right unitor with a composite. For successive
restriction, the pentagon compares the associators on the two routes.
Reassociate vertical composites to obtain the required matching equation.
\end{proof}

\begin{lemma}
\label[lemma]{lem:Supplementary_Cone_Flattening}
Let $f\colon A\to B$, $g\colon B\to Z$, $h\colon X\to Z$, and let
$t=(u,v,\nu)$ be a cone over $(g,h)$ on $Q$.
A cone $s=(p,q,\rho)$ over $(f,u)$ on $T$ determines a cone
$\operatorname{Flat}_t(s)$ over $(gf,h)$ with legs $p,vq$ and matching
\[
 \operatorname{match}(q^*t)\cdot(g\star\rho)\cdot\alpha_{p,f,g}.
\]
Flattening takes cone comparisons to cone comparisons. Moreover,
$r^*\operatorname{Flat}_t(s)$ compares with
$\operatorname{Flat}_t(r^*s)$ by the identity on the left leg and
$\alpha_{r,q,v}$ on the right leg.
\end{lemma}
\begin{proof}[Proof idea]
Postwhisker the matching equation for a comparison of the inner cones by
$g$, and paste it with the equation for the corresponding restrictions
of $t$. The resulting equation is precisely compatibility of the
flattened comparisons.

For restriction, the intermediate cone obtained from $t$ can be restricted
either successively by $q,r$, or by $qr$.
Apply \Cref{lem:Supplementary_Cone_Restriction} to compare these cones,
then flatten. The leg comparisons normalize to the identity and the
displayed associator. Their matching equation is retained during this normalization.
\end{proof}

\begin{lemma}
\label[lemma]{lem:Supplementary_Pasting_Factorization}
Suppose $t$ is a pullback cone. Put $N=A\times_B Q$ and
$O=A\times_Z X$. Flattening the defining cone of $N$ induces an equivalence
$F\colon N\to O$. If $\ell_s\colon T\to N$ and
$\ell_{\operatorname{Flat}_t(s)}\colon T\to O$ are the chosen factorization
functors, then
\[
 \ell_{\operatorname{Flat}_t(s)}=F\circ\ell_s.
\]
Consequently, $s$ is a pullback cone if and only if
$\operatorname{Flat}_t(s)$ is a pullback cone.
\end{lemma}
\begin{proof}[Proof idea]
The defining cone of $O$ gives a cone over $(g,h)$. Factor it through $t$
to obtain $O\to Q$, then lift the resulting cone over $(f,u)$ to obtain
$O\to N$. The two inverse comparisons follow from the pullback universal
properties, using full cone comparisons.

For the displayed factorization, flatten the computation comparison for
$\ell_s$ and apply \Cref{lem:Supplementary_Cone_Flattening} to its restriction.
This compares the cone induced by $F\ell_s$ with the flattened input cone.
Reflect through the pullback property of $O$. A cone is a pullback exactly
when its factorization into the chosen pullback is an equivalence;
cancellation with $F$ proves the final assertion.
\end{proof}
\end{document}
